\section{Technical reference}\label{app:tech}

\subsection*{Opacity and equation-of-state continuations}
Settling raises the core metal fraction above the $Z=0.03$ limit of the
tabulated ATOMIC opacities. Through $Z=0.08$ the radiative opacity is
continued linearly from the $Z=0.02$ and 0.03 sources at fixed temperature,
density and hydrogen abundance,
\begin{equation}
 \kappa(T,\rho,X,Z)=\kappa(T,\rho,X,0.03)
 +\frac{Z-0.03}{0.01}\,[\kappa(T,\rho,X,0.03)-\kappa(T,\rho,X,0.02)],
\end{equation}
with all source temperature and density bounds enforced. Transferring the
OPLIB composition ratio from $Z=0.02$ to a withheld $Z=0.03$ source differs
from the ATOMIC value by at most 0.1867\% at 34 core states, and the linear
and ratio forms differ by 0.001212\% in luminosity over a matched 5 Gyr.
Independent OPLIB data through $Z=0.2$ support linear scaling within
3.302\% at 6--10 MK; in the region with $Z>0.06$ radiation carries at most
0.04154\% of the gradient-driven heat. As the envelope falls below
$Z=0.01$ the molecular opacity interpolates the actual AESOPUS values,
including $Z=0$ and 0.004; at higher temperatures $\kappa$ is continued
linearly in $Z$ from the 0.01 and 0.02 sources with a floor at
$Z=10^{-4}$. Hydrogen dependence is continued in $\log\kappa$ from the
$X=0.70$ and 0.75 sources through pure hydrogen; the molecular source ends
at $1-X=10^{-7}$. In the dense hydrogen envelope, where the OPLIB table
ends at $\log R=1.5$ ($R=\rho/T_6^3$), the TOPS hydrogen slope is blended
in over $\log R=1.25$--1.45 and used to $\log R\leq1.8$ for
$\log T=5.6$--6.1; convection carries 99.96\% of the heat at the affected
layer. In dense core cells without spectral coverage (0.125--0.45 keV,
$10^4$--$10^5$ g cm$^{-3}$, $X\leq0.2$) the tabulated uncut opacity is
multiplied by an inferred plasma correction; radiation carries at most
0.2535\% of the local transport there.

The variable-metal equation-of-state family covers $Z=0$--0.08 (0.16 for
the core) and all 23 sampled Track S structures. Against direct FreeEOS
evaluations the maximum differences are 0.03028\% in pressure, 0.09056\% in
adiabatic gradient, 0.7642\% in heat capacity and 1.037\% in $\Gamma_1$;
36 low-density and 28 high-metal comparisons agree within 0.009197\% and
0.02530\% in the material responses. The Track F family (152 tables, 38
hydrogen and four $^3$He fractions) has a maximum second-composition-derivative
difference of 0.1499\% at 102 compositions and a heat-capacity difference of
0.1062\% on the saved 512-point structure. Composition derivatives are
evaluated without an abundance floor, including when a species is exactly
absent; 128 analytic comparisons and five FreeEOS curvature checks at zero
$^3$He (within 0.007181\%) pass. The continuous driver's potential matches
direct FreeEOS to 0.02589\% in pressure and 0.02467\% in internal energy on
the 512-zone reference profile.

\subsection*{Atmosphere families}
Table~\ref{tab:atm} lists the atmosphere tables entered by Track S and the
flash continuations. ``Depth'' is the largest change in joining temperature
and pressure when the source column is computed to a deeper lower boundary;
``interpolation'' is the largest difference between an independently solved
column and the interpolated value. Missing cells remain inaccessible to the
evolution rather than being filled.

{\small\renewcommand{\arraystretch}{1.1}
\begin{longtable}{@{}>{\raggedright\arraybackslash}p{0.30\textwidth}>{\raggedright\arraybackslash}p{0.22\textwidth}rp{0.12\textwidth}p{0.16\textwidth}@{}}
\caption{Atmosphere tables along Track S and the flash continuations.}\label{tab:atm}\\
\toprule Composition & $\Teff$ (K), $\log g$ & Columns & Depth ($T$; $P$) & Interpolation ($T$; $P$) \\\midrule\endfirsthead
\toprule Composition & $\Teff$ (K), $\log g$ & Columns & Depth ($T$; $P$) & Interpolation ($T$; $P$) \\\midrule\endhead
Track P reference: $X=0.1$--0.7, $X_3=0$--0.12, $Z=0.02$ & to 5600; 4.9--5.4 & 220 & 0.007883\% & 0.7399\%; 0.5707\% \\
Warm extension, $X=0.15$, 0.2 & 5400--6400; 5.7, 6.0 & 300 & 0.009372\% & 0.4860\% \\
Late Track F grid, $X=0.70$--0.90, $X_3=0$--0.03 & 3400--4000; 4.9--5.15 & 48 & --- & 1.837\% in $T$ \\
Metal response $Z=0.01$--0.02 (paired) & to 4400; 5.025 & pairs & 0.0002576\% & 0.6093\%; 3.064\% \\
$Z=0.005$--0.01 & to 4400; 5.1 & pairs & --- & 1.671\%; 1.980\% \\
$Z=0.0025$--0.005 & to 4400; 5.1 & 12 pairs & 0.001301\%; 0.002243\% & 0.03091\%; 0.2311\% \\
$Z=0.001$--0.0025 & 4000--4400; 5.1 & 8 pairs & --- & 0.2539\%; 3.111\% \\
$Z=0.0005$--0.001 & 4000--4400; 5.1 & 8 pairs & --- & 2.031\%; 2.755\% \\
$X=0.9$--0.98, $Z=0.001$, $X_3=0$, 0.003 & 4000--4800; 4.9--5.4 & 70 & --- & 0.004095\%; 0.01689\% \\
$X=0.995$--0.9995, $Z=10^{-6}$--$10^{-4}$ & 4600--5200; 5.15--5.65 & 94 & --- & 0.09719\%; 4.656\% \\
$1-X=10^{-5}$--$10^{-3}$, $Z=10^{-8}$--$10^{-6}$ & 4600--5200; 5.4--5.65 & 81 & (shallowest ends at $\tau=168.3$) & 0.4624\%; 0.3930\% \\
$Z\leq10^{-6}$, He $10^{-10}$ & 4700--6000; 5.5--5.95 & 72 & $4.327\times10^{-6}$\%; $4.770\times10^{-5}$\% & 0.1751\%; 0.5986\% \\
Zero metal, cold & 4650--4800; 5.9--5.95 & 6 & below $10^{-6}$ & 0.1154\%; 1.108\% \\
Gravity continuation & 4650--4800; 5.95--6.1 & inferred & --- & 0.1878\%; 2.391\% (9 solved) \\
Temperature continuation & 4600; 5.9 & 1 solved & --- & within 0.05\% \\
H--He mixtures (flash), $X=0.78$--0.9955 & 4500--5500; 4.95--5.80 & 54 & 0.02452\%; 0.03703\% & 0.08784\%; 1.094\% \\
H--He mixtures, three $X$ & 4700--6000; 4.95--5.25 & 36 & 0.01078\%; 0.06911\% & --- \\
H--He mixtures, three $X$ & 6000--6500; 4.20--5.55 & 60 & 0.01373\%; 0.09599\% & --- \\
H--He mixtures, three $X$ & 5000--6000; 5.4--5.8 & 27 & --- & $7.587\times10^{-6}$\%; 0.1582\% \\
MESA hydrogen white-dwarf table \cite{mesa_atm} & 2000--40,000; 5.5--9.5 & table & joined at $\tau=25.12$ & extrapolated below $\log g=5.5$ \\
\bottomrule
\end{longtable}}

Within $0.90<X<0.94$ the logarithms of the joining values are interpolated
in hydrogen between the $X=0.90$ and 0.94 families; the metal response is
calibrated at $X_3=0$ and applied separately from the small isotope
response, and helium below a mass fraction 0.005 is neglected in the
4600--4800 K, $\log g=5.9$--6.1 domain (solved helium-bearing columns change
the joining values by less than 0.1\% in $T$ and 2.5\% in $P$). In the flash
continuations the hydrogen table is extrapolated below $\log g=5.5$ using
either its lowest two rows or the gravity response of solved near-hydrogen
atmospheres at $\log g=4.35$ and 4.50; the two inferred references differ by
up to 2.189\% in temperature and 15.04\% in pressure at tabulated points,
and paired stellar comparisons between them are those quoted in
Section~\ref{sec:flash}.

\subsection*{Time-step and conservation checks}
Every accepted interval compares one full step with two half steps in
structure, species, total hydrogen and energy. Species balances are checked
independently in each composition region together with the inert-metal
mass, nuclear rest-mass release and the first law. For Tracks S, F and P, the relative time-error
budgets are 0.0003 for radius, density and temperature, 0.003 for local
$^3$He and 0.001 for the other species; the abundance solver tolerance is
$10^{-10}$ with a stricter fallback (a $10^{-12}$ control changes mass
fractions by at most $1.360\times10^{-10}$). During the flash, deposited
nuclear energy is integrated consistently with the implicit equations, time
refinement permits 0.5\% in deposited energy and 0.1\% in structure and
endpoint nuclear power (a 1\% allowance was tested and rejected because a
local temperature difference of $|\Delta\ln T|=0.001091$ exceeded the
declared 0.001 limit), and the first-law residual is measured against the
larger of nuclear and photon luminosity with a limit of $2\times10^{-8}$.
Mesh insertions preserve species inventories, internal energy and the nodal
gravitational-energy sum. Increasing the separate abundance-change limit per
step from 0.001 to 0.005 saves 24.99\% of CPU over a matched 1 Gyr with a
luminosity change of 0.0005727\%; a larger structure tolerance was rejected
(Section~\ref{sec:model}). Electron-pair collision integrals for the late
shell-burning calculation cover degeneracy parameter $\eta=-6$--64 and
screening parameter $b=0.1$--2400, with interpolation and integration checks
within 0.3466\%. Collision integrals use 96 Jacobi and 128 Legendre points;
against 384/512 points the maximum scaled coefficient difference is
$4.851\times10^{-9}$.
